\documentclass[leqno]{article}
% Shared QBT/Soln pack preamble for resources.danieljsmith.org.
\usepackage{graphicx}
\usepackage{dsfont}
\usepackage{mathtools}
\usepackage{amssymb}
\usepackage{amsmath}
\usepackage{mleftright}
\usepackage{amsthm}
\usepackage{float}
\usepackage{ragged2e}
\usepackage{tensor}
\usepackage{fancyhdr}
\usepackage{empheq}
\usepackage[most]{tcolorbox}
\usepackage{enumitem}
\usepackage{dirtytalk}
\usepackage{marginnote}
\usepackage{microtype}
\usepackage[
  a4paper,
  left=0.8in,
  right=1in,
  top=1in,
  bottom=0.5in,
  headheight=14pt,
  headsep=0.4in
]{geometry}
\setlength{\parindent}{0cm}

% TikZ / PGFPlots
\usepackage{pgfplots}
\pgfplotsset{compat=1.18}
\usepgfplotslibrary{fillbetween, polar}
\usetikzlibrary{calc, positioning}
\usetikzlibrary{angles, quotes}
\usetikzlibrary{arrows.meta}
\usetikzlibrary{decorations.pathreplacing, decorations.markings}
\usetikzlibrary{patterns}
\usetikzlibrary{intersections}

% Links and contents. The contents follow the TMUA papers: compact entries,
% no dot leaders, and each question's mark total in grey after its label.
\usepackage{tocloft}
\tocloftpagestyle{fancy}
\setlength{\cftbeforesubsecskip}{2pt}
\setlength{\cftsubsecindent}{1.5em}
\setlength{\cftsubsecnumwidth}{0pt}
\renewcommand{\cftsubsecfont}{\small\color{black}}
\renewcommand{\cftsubsecpagefont}{\small\color{black}}
\renewcommand{\cftsubsecleader}{\hfill}
\usepackage[colorlinks=true,linkcolor=black,citecolor=magenta,urlcolor=black]{hyperref}%
\urlstyle{same}

% Pack macros
% \marks also adds its value to the current question's total in the .aux
% file, so the contents show each total from the second pass. The write
% happens at shipout, so a box typeset more than once counts once. The
% pack's mark total for the contents-page title is summed the same way.
\makeatletter
\newcount\djs@totalmarks
\newcommand{\djsAddMarks}[2]{%
  \global\advance\djs@totalmarks#2\relax
  \@ifundefined{djs@marks@#1}%
    {\expandafter\xdef\csname djs@marks@#1\endcsname{#2}}%
    {\expandafter\xdef\csname djs@marks@#1\endcsname
      {\the\numexpr\csname djs@marks@#1\endcsname+#2\relax}}}
\RenewDocumentCommand{\marks}{o m}{%
  \protected@write\@auxout{}{\string\djsAddMarks{\number\value{enumi}}{#2}}%
  \IfValueTF{#1}%
    {\marginnote{\raggedright\textbf{[#2]}}[#1]}%
    {\marginnote{\raggedright\textbf{[#2]}}}%
}
\makeatother

\newcommand{\vspacerule}[2]{%
  \par\vspace*{#1}%
  \noindent\hrulefill\par
  \vspace*{#2}%
}

\definecolor{scarlet}{RGB}{205, 40, 75}
% Site colours for the contents-page title: the accent, the rule colour, a
% stronger rule colour for the path separator, and the muted text colour.
\definecolor{djsBurgundy}{RGB}{118, 59, 67}
\definecolor{djsHairline}{RGB}{216, 213, 205}
\definecolor{djsSlash}{RGB}{170, 166, 158}
\definecolor{djsMuted}{RGB}{96, 101, 106}

% Worked-solution box, adapted from the TMUA worked-solutions box. A white breakable box with a left scarlet rule and no surrounding
% frame or tint. A part that breaks fills the rest of its page, so the rule
% runs to the bottom margin; the unbroken or last part ends in a short
% horizontal foot rule so a solution visibly stops. Unlike the TMUA box there
% is no answer label. The box does not grow to the left, so the rule sits at
% the text edge; growing to the right keeps the usable line width close to the
% old tcolorbox Solution box.
% Spacing inside the box: paragraphs are 6pt apart, so source blank lines show
% (solutions put one after each display and none between a lead-in and its
% display). Every display gets 6pt above and 10pt below, whatever the length
% of its lead-in line; align rows are spread by an 8pt \jot; solution parts
% are 14pt apart and sub-parts 8pt; \left and \right space like ordinary
% brackets (mleftright); no single line is stranded at a page break; and
% neither the Solution label nor the line before a list ends a page. No
% skip can shrink, so full pages are not squeezed.
\newtcolorbox{djsSolution}{%
  enhanced, breakable, frame hidden, sharp corners, colback=white,
  height fixed for=first and middle, height fill,
  extras unbroken and last={natural height},
  borderline west={2.2pt}{0pt}{scarlet},
  left=11pt, right=0pt, top=5pt, bottom=13pt,
  before skip=13pt, after skip=4pt,
  grow to right by=10mm,
  overlay unbroken and last={%
    \fill[scarlet] (frame.south west)
      rectangle ([xshift=3.5cm, yshift=2.2pt]frame.south west);},
  before upper={{\color{scarlet}\bfseries Solution}\par\nobreak\vspace{4pt}%
    \setlength{\parskip}{6pt plus 2pt}%
    \setlength{\abovedisplayskip}{6pt plus 2pt}%
    \setlength{\abovedisplayshortskip}{6pt plus 2pt}%
    \setlength{\belowdisplayskip}{10pt plus 2pt}%
    \setlength{\belowdisplayshortskip}{10pt plus 2pt}%
    \setlength{\jot}{8pt}%
    \setlist{parsep=6pt plus 2pt, topsep=6pt plus 2pt, beginpenalty=10000}%
    \setlist[enumerate,1]{itemsep=14pt plus 2pt}%
    \setlist[enumerate,2]{itemsep=8pt plus 2pt}%
    \mleftright
    \clubpenalty=10000 \widowpenalty=10000 \displaywidowpenalty=10000\relax}}

% Optional smaller figure inside a solution box. \scalebox shrinks node text
% with the drawing. Not applied to existing figures.
\newcommand{\SolnFigureScale}{0.85}
\newsavebox{\djsSolnFigureBox}
\newenvironment{djsSolutionFigure}
  {\par\vspace{4pt}\begin{center}\begin{lrbox}{\djsSolnFigureBox}}
  {\end{lrbox}\scalebox{\SolnFigureScale}{\usebox{\djsSolnFigureBox}}%
   \end{center}\vspace{2pt}\par}

\newcommand{\cxl}[2]{%
  \tikz[baseline=(X.base)]{%
    \node[inner sep=0pt, outer sep=0pt, text=#1] (X) {$#2$};
    \draw[#1, line width=0.4pt] (X.south west)--(X.north east);
  }%
}

% Page-style hook run at the contents page and each question's opening page.
% A no-op for QBT packs; \djsSolnHeader points it at the fancy header.
\newcommand{\djsQuestionPageStyle}{}

% A contents entry: "Question N" in a box wide enough for two digits, so the
% mark totals line up, then the total in grey. Bookmarks keep "Question N".
\makeatletter
\newlength{\djsTocLabelWidth}
\newcommand{\djsTocMarks}[1]{%
  \@ifundefined{djs@marks@#1}{}{%
    {\footnotesize\itshape\color{black!45}%
     \csname djs@marks@#1\endcsname~mark%
     \ifnum\csname djs@marks@#1\endcsname=1 \else s\fi}}}
\newcommand{\djsTocEntry}[1]{%
  \texorpdfstring
    {\settowidth{\djsTocLabelWidth}{Question 00}%
     \makebox[\dimexpr\djsTocLabelWidth+0.9em\relax][l]{Question #1}%
     \djsTocMarks{#1}}%
    {Question #1}}
\makeatother

% The question count for the contents-page title is written to the .aux at
% the end of the document, so it also appears from the second pass.
\newcounter{djsQuestions}
\newcommand{\questionitem}{%
  \djsQuestionPageStyle
  \item\phantomsection
  \stepcounter{djsQuestions}%
  \addcontentsline{toc}{subsection}{\protect\djsTocEntry{\number\value{enumi}}}%
}
\makeatletter
\newcommand{\djsSetQuestionCount}[1]{\gdef\djs@questioncount{#1}}
\AtEndDocument{\immediate\write\@auxout
  {\string\djsSetQuestionCount{\number\value{djsQuestions}}}}
\makeatother

\renewcommand{\contentsname}{Questions}

% \djsPackHeader[<link content tag>]{<left header>}; the tag says which
% document a click on the header link came from.
\newcommand{\djsPackHeader}[2][questions]{%
  \pagestyle{fancy}%
  \fancyhead{}%
  \fancyfoot{}%
  \renewcommand{\headrulewidth}{0.9pt}%
  \lhead{#2}%
  \chead{}%
  \rhead{\href{https://danieljsmith.org/?utm_source=pdf&utm_campaign=further-maths&utm_content=#1}{danieljsmith.org}}%
}

% The header's topic is also the contents-page title, under which a label
% says which booklet this is.
\newcommand{\djsPackTitle}{}
\newcommand{\djsBookletLabel}{}
\newcommand{\djsQbtHeader}[1]{%
  \djsPackHeader{#1}%
  \renewcommand{\djsPackTitle}{#1}%
  \renewcommand{\djsBookletLabel}{Questions by topic}%
}

% The colophon at the foot of page one, independent of the page style: the
% contents-page title's accent rule (the same \djs@accentrule, so the two
% always match), then the date on the left and a link to the site on the
% right, all inside the bottom margin. Every later page carries its number
% over the page count at the bottom right, under the header's link. The count
% is the kernel's \PreviousTotalPages, which every run already records in the
% .aux file. Sizes are fixed, so these footers match the TMUA papers'. The
% document supplies the date; recompiling never changes it.
\makeatletter
\newcommand{\djsLastUpdated}[1]{%
  \AtBeginDocument{\AddToHookNext{shipout/foreground}{\djs@colophon{#1}}}}
\newcommand{\djs@footR}[1]{%
  \put(\dimexpr\paperwidth-1in\relax,-\dimexpr\paperheight-0.32in\relax){%
    \makebox[0pt][r]{#1}}}
\newcommand{\djs@colophon}[1]{%
  \put(0.8in,-\dimexpr\paperheight-0.47in\relax){%
    \makebox[0pt][l]{\djs@accentrule{\textwidth}}}%
  \put(0.8in,-\dimexpr\paperheight-0.32in\relax){\makebox[0pt][l]{%
    \fontsize{9}{11}\selectfont Last updated #1}}%
  \djs@footR{\fontsize{7}{8.5}\selectfont\color{djsMuted}%
    Latest version at \href{https://resources.danieljsmith.org/}%
      {\color{black}resources.danieljsmith.org}}}
\AddToHook{shipout/foreground}{%
  \ifnum\value{page}>1
    \djs@footR{{\fontsize{9}{11}\selectfont\thepage}%
      {\fontsize{7}{8.5}\selectfont\color{djsMuted}%
        \hspace{0.3em}/\hspace{0.3em}\PreviousTotalPages}}%
  \fi}
\makeatother

% Solutions packs show the header only on the contents page and each
% question's opening page, so it never appears part-way through a solution.
% \thispagestyle marks the next page shipped out, which is the question's
% page because every solution question starts after a \newpage.
\newcommand{\djsSolnHeader}[1]{%
  \djsPackHeader[solutions]{\textit{Solutions} - #1}%
  \pagestyle{empty}%
  \renewcommand{\djsQuestionPageStyle}{\thispagestyle{fancy}}%
  \renewcommand{\djsPackTitle}{#1}%
  \renewcommand{\djsBookletLabel}{Worked solutions}%
}

% \djsSitePath{<strand>}{<topic>}, called once after the header, names where
% the pack sits on the site; the contents-page title prints it above the pack
% title, e.g. CORE PURE / FURTHER CALCULUS. Without it the line is left out.
\makeatletter
\newif\ifdjs@sitepath
\newcommand{\djs@strand}{}
\newcommand{\djs@section}{}
\newcommand{\djsSitePath}[2]{%
  \djs@sitepathtrue
  \renewcommand{\djs@strand}{#1}%
  \renewcommand{\djs@section}{#2}%
}
\makeatother

% \djsBoards{\djsBoard{<label>}{<code>}...}, called once after \djsSitePath,
% lists the pack's exam boards (every board in its strand for a pack for all
% boards); the contents-page title prints them under the booklet label, after
% a burgundy EXAM BOARDS key. A board's label and code never break across
% lines: a line break falls only between boards, where it replaces the
% separating dot.
\makeatletter
\newif\ifdjs@boards
\newif\ifdjs@firstboard
\newcommand{\djs@boardlist}{}
\newcommand{\djsBoards}[1]{%
  \djs@boardstrue
  \renewcommand{\djs@boardlist}{#1}%
}
\newcommand{\djsBoard}[2]{%
  \ifdjs@firstboard\global\djs@firstboardfalse\else
    \discretionary{}{}{\kern0.45em\textperiodcentered\kern0.45em}%
  \fi
  \mbox{#1 {\color{djsMuted}#2}}}
\newlength{\djs@boardkeywd}
\newcommand{\djs@boardkey}{{\color{djsBurgundy}\scriptsize\bfseries\djs@caps{Exam boards}}}
\newcommand{\djs@boardsline}{%
  \ifdjs@boards
    \vspace{16pt}%
    \settowidth{\djs@boardkeywd}{\djs@boardkey}%
    \noindent\makebox[\dimexpr\djs@boardkeywd+1.1em\relax][l]{\djs@boardkey}%
    \parbox[t]{\dimexpr\linewidth-\djs@boardkeywd-1.1em\relax}{%
      \raggedright\small\global\djs@firstboardtrue\djs@boardlist}\par
  \fi}
\makeatother

% Further Mechanics packs call \djsMechanicsPack in their preamble, so that
% \djsFrontMatter prints the default value of g above the contents. The box
% keeps its own space above it (before skip is the whole gap below the title,
% not an addition to it). Without the box, \djsGravityNote is just extra space
% before the contents heading, bringing the gap above it close to the gap
% between the box and the heading on Further Mechanics pages.
\newcommand{\djsGravityNote}{\vspace{14pt}}
\newcommand{\djsMechanicsPack}{%
  \renewcommand{\djsGravityNote}{%
    \begin{tcolorbox}[sharp corners, colback=white, colframe=black,
      boxrule=1pt, left=8pt, right=8pt, top=6pt, bottom=6pt,
      before skip=20pt]
    {\bfseries Instructions}\par\vspace{4pt}
    Unless a question states otherwise, take $g = 9.8\,\mathrm{m\,s^{-2}}$
    wherever a numerical value of $g$ is required. Solutions also give final
    answers for $g = 10\,\mathrm{m\,s^{-2}}$ where these differ.
    \end{tcolorbox}%
    \vspace{16pt}}%
}

% Contents-page title: the site path in spaced burgundy capitals, the pack
% title, a hairline led by a short scarlet rule (the solution box's foot rule),
% then the booklet label with the question and mark counts, and the exam
% boards line when the pack has one. The counts come from the .aux, so they
% appear from the second pass.
\makeatletter
\newsavebox{\djs@eyebrowbox}
\newcommand{\djs@caps}[1]{\textls[180]{\MakeUppercase{#1}}}
\newcommand{\djs@eyebrow}{%
  \color{djsBurgundy}\footnotesize\bfseries
  \djs@caps{\djs@strand}%
  {\color{djsSlash}\hspace{0.55em}/\hspace{0.55em}}%
  \djs@caps{\djs@section}}
\newcommand{\djs@dot}{\hspace{0.45em}\textperiodcentered\hspace{0.45em}}
\newcommand{\djs@counts}{%
  \@ifundefined{djs@questioncount}{}{%
    \djs@questioncount~question\ifnum\djs@questioncount=1 \else s\fi
    \ifnum\djs@totalmarks>0
      \djs@dot\the\djs@totalmarks~mark\ifnum\djs@totalmarks=1 \else s\fi
    \fi}}
% The accent rule: a hairline of the given width, led by a short scarlet rule
% (the solution box's foot rule). The contents-page title and the colophon
% both draw it, so their weights are always the same.
\newcommand{\djs@accentrule}[1]{%
  \makebox[0pt][l]{\color{djsHairline}\rule[0.9pt]{#1}{0.6pt}}%
  {\color{scarlet}\rule{3.5cm}{2.2pt}}}
\newcommand{\djsContentsTitle}{%
  \ifdjs@sitepath
    \sbox\djs@eyebrowbox{\djs@eyebrow}%
    \ifdim\wd\djs@eyebrowbox>\linewidth
      \@latex@warning@no@line{The \string\djsSitePath\space line is wider
        than the text, so it wraps. Shorten its topic}%
    \fi
    {\raggedright\djs@eyebrow\par}%
    \vspace{7pt}%
  \fi
  {\raggedright\fontsize{24.88}{28}\selectfont\djsPackTitle\par}%
  \vspace{9pt}%
  \noindent\djs@accentrule{\linewidth}\par
  \vspace{5pt}%
  {\small\color{djsMuted}\djsBookletLabel\hfill\djs@counts\par}%
  \djs@boardsline
  \vspace{8pt}%
}
\makeatother

\newcommand{\djsFrontMatter}{%
  \djsQuestionPageStyle
  \vspace*{6pt}%
  \djsContentsTitle
  \djsGravityNote
  \tableofcontents
  \newpage
}


\djsQbtHeader{Simple Harmonic Motion}
\djsSitePath{Further Mechanics}{Circular Motion \& Oscillations}
\djsBoards{\djsBoard{Edexcel FM2}{9FM0/4C}\djsBoard{OCR MEI Mech Major}{Y421}}
\djsMechanicsPack
\djsLastUpdated{6 November 2026}

\begin{document}
\djsFrontMatter
\begin{enumerate}[
  label=\textbf{\arabic*.},
  leftmargin=*,
  topsep=0pt,
  partopsep=0pt,
  parsep=0pt
]

% ---- Question 1 ----
\questionitem
% [Source: _QBT__Simple_Harmonic_Motion, Q1]

A particle $P$ is moving along the $x$-axis. At time $t$ seconds, the displacement of $P$ from the origin $O$ is $x$ metres, where
\[
x = 3\cos\left(\tfrac{1}{4}\pi t\right) + 3\sqrt{3}\sin\left(\tfrac{1}{4}\pi t\right)
\]
\\
\begin{enumerate}[label=\textbf{(\alph*)}]
\item Prove that $P$ is moving with simple harmonic motion. \marks{3} \\
\item Find the period of the motion. \marks{2} \\
\item Find the amplitude of the motion. \marks{2} \\
\item Find, in terms of $\pi$, the maximum speed of $P$. \marks{2} \\
\item Find the total time during each complete oscillation for which the speed of $P$ is greater than $\frac{3\sqrt{3}}{4}\pi\,\mathrm{m\,s^{-1}}$. \marks{6} \\
\end{enumerate}

\newpage

% ---- Question 2 ----
\questionitem
% [Source: _QBT__Simple_Harmonic_Motion, Q2]

A particle $P$ of mass $0.5\,\mathrm{kg}$ moves in a straight line with simple harmonic motion about a centre $O$. The time taken for one complete oscillation is $\frac{\pi}{2}$ seconds. \\
When $P$ is $0.3\,\mathrm{m}$ from $O$, its speed is $1.6\,\mathrm{m\,s^{-1}}$. \\
\begin{enumerate}[label=\textbf{(\alph*)}]
\item Show that the amplitude of the motion is $0.5\,\mathrm{m}$. \marks{4} \\
\item Find the length of time within each complete oscillation for which the speed of $P$ is greater than $1\,\mathrm{m\,s^{-1}}$, giving your answer as an exact multiple of $\pi$. \marks{6} \\
\item Find the magnitude of the resultant force acting on $P$ at an instant when the kinetic energy of $P$ is $0.36\,\mathrm{J}$. \marks{4} \\
\end{enumerate}

\newpage

% ---- Question 3 ----
\questionitem
% [Source: _QBT__Simple_Harmonic_Motion, Q3]

A particle $P$ moves in a straight line with simple harmonic motion about a fixed point $O$. \\
When $P$ is $0.5\,\mathrm{m}$ from $O$, its speed is $3.6\,\mathrm{m\,s^{-1}}$, and when $P$ is $1.2\,\mathrm{m}$ from $O$, its speed is $1.5\,\mathrm{m\,s^{-1}}$. \\
The amplitude of the motion is $a$ metres. \\
\begin{enumerate}[label=\textbf{(\alph*)}]
\item Show that $a=1.3$. \marks{5} \\
\item Find the exact period of the motion and the greatest magnitude of the acceleration of $P$. \marks{3} \\
\end{enumerate}
During one complete oscillation, the magnitude of the acceleration of $P$ is at least $6.3\,\mathrm{m\,s^{-2}}$ for a total of $T$ seconds. \\
\begin{enumerate}[resume, label=\textbf{(\alph*)}]
\item Find the value of $T$. \marks{6} \\
\end{enumerate}

\newpage

% ---- Question 4 ----
\questionitem
% [Source: _QBT__Simple_Harmonic_Motion, Q4]

The depth of water at the entrance to a harbour varies with the tides, and the water level moves with simple harmonic motion. \\
On a particular day, high tide occurs at 05:40, when the depth of water is $11.2\,\mathrm{m}$, and the next low tide occurs at 11:55, when the depth of water is $3.6\,\mathrm{m}$. \\
\begin{enumerate}[label=\textbf{(\alph*)}]
  \item Between 05:40 and 11:55 there are two moments at which the water level is falling at a rate of $1.5\,\mathrm{m\,h^{-1}}$. Find the depth of water at each of these moments. \marks{6} \\
\end{enumerate}
A ship can only use the harbour when the depth of water at the entrance is at least $5.5\,\mathrm{m}$. \\
\begin{enumerate}[label=\textbf{(\alph*)}, resume]
  \item Find the length of time between 05:40 and the following high tide during which the ship cannot use the harbour. \marks{4} \\
\end{enumerate}

\newpage

% ---- Question 5 ----
\questionitem
% [Source: _QBT__Simple_Harmonic_Motion, Q5]

A particle $P$ moves along the $x$-axis with simple harmonic motion. At time $t$ seconds, its displacement from the origin $O$ is $x$ metres, where
\[
\frac{\mathrm{d}^2x}{\mathrm{d}t^2}=9(2-x)
\]
and $P$ is instantaneously at rest when $x=5$. \\
\begin{enumerate}[label=\textbf{(\alph*)}]
\item Show that $v$, the velocity of $P$ in $\mathrm{m\,s^{-1}}$ at time $t$, satisfies
\[
v^2=9(5-x)(x+1)
\]
\marks[-20pt]{5}
\item When $t=0$, $P$ is at the point where $x=\frac{1}{2}$ and is moving in the positive $x$-direction. \\
Hence, or otherwise, show that
\[
x=2+3\sin\left(3t-\frac{\pi}{6}\right)
\]
\marks[-30pt]{6}
\end{enumerate}

\newpage

% ---- Question 6 ----
\questionitem
% [Source: _QBT__Simple_Harmonic_Motion, Q6]

A box of mass $m\,\mathrm{kg}$ rests on a horizontal platform. The platform moves vertically with simple harmonic motion of angular frequency $\omega\,\mathrm{rad\,s^{-1}}$ and amplitude $a\,\mathrm{m}$ about a mean position $O$. At time $t$ seconds, the upward displacement of the platform from $O$ is $x\,\mathrm{m}$. \\
The box remains in contact with the platform throughout the motion, and the greatest magnitude of the force exerted by the platform on the box is $\frac{7}{3}$ times the least magnitude of this force. \\
\begin{enumerate}[label=\textbf{(\alph*)}]
\item Show that $\omega^2=\frac{2g}{5a}$. \marks{3} \\
\end{enumerate}
The angular frequency is unchanged, but the amplitude of the motion is increased to $3a\,\mathrm{m}$. \\
\begin{enumerate}[label=\textbf{(\alph*)}, resume]
\item Show that the box leaves the platform, and find, in terms of $a$, the height above $O$ at which this first happens. \marks{3} \\
\item Find, in terms of $a$ and $g$, the speed of the box at the instant it leaves the platform. \marks{3} \\
\end{enumerate}

\newpage

% ---- Question 7 ----
\questionitem
% [Source: _QBT__Simple_Harmonic_Motion, Q7]

Throughout this question, use $g = 10\,\mathrm{m\,s^{-2}}$. \\
A light elastic string has natural length $0.9\,\mathrm{m}$ and modulus of elasticity $36\,\mathrm{N}$. \\
A particle $P$ of mass $0.4\,\mathrm{kg}$ is attached to one end of the string. The other end of the string is attached to a fixed point $A$. Particle $P$ hangs freely in equilibrium with $P$ vertically below $A$. \\
The particle is then pulled vertically down to a point $B$ and released from rest. \\
\begin{enumerate}[label=\textbf{(\alph*)}]
\item Show that, while the string is taut, $P$ moves with simple harmonic motion with period $\frac{\pi}{5}$ seconds. \marks{6} \\
\end{enumerate}
When the string first becomes slack, the speed of $P$ is $\sqrt{3}\,\mathrm{m\,s^{-1}}$. \\
\begin{enumerate}[resume, label=\textbf{(\alph*)}]
\item Find the distance $AB$. \marks{3} \\
\end{enumerate}
After the string becomes slack, $P$ continues to move vertically upwards until it first comes to instantaneous rest at the point $C$. \\
\begin{enumerate}[resume, label=\textbf{(\alph*)}]
\item Find the distance $BC$ and the exact time taken for $P$ to move from $B$ to $C$. \marks{6} \\
\end{enumerate}

\newpage

% ---- Question 8 ----
\questionitem
% [Source: _QBT__Simple_Harmonic_Motion, Q8]

\begin{figure}[H]
    \centering
\begin{tikzpicture}[scale=1.3]
\pgfmathsetmacro{\ang}{asin(3/5)}
\def\planelen{7.2}
\def\groundlen{6.4}
\def\posA{0.42}
\def\posO{0.9}
\def\gap{0.07}
\def\prad{1.8pt}

\coordinate (base) at (0,0);
\coordinate (top) at (\ang:\planelen);
\coordinate (horiz) at (\groundlen,0);
\coordinate (A) at ($(base)!\posA!(top)$);
\coordinate (O) at ($(base)!\posO!(top)$);
\coordinate (P) at ($(A)+({\ang+90}:\gap)$);
\coordinate (Os) at ($(O)+({\ang+90}:\gap)$);

\draw[thick] (base) -- (horiz);
\draw[thick] (base) -- (top);

\draw[thick] (P) -- (Os) -- (O);

\fill (P) circle (\prad);

\pic["$\alpha$", draw, angle radius=10mm, angle eccentricity=1.35] {angle=horiz--base--top};

\node[above left] at (P) {$P$};
\node[below right] at (A) {$A$};
\node[below right] at (O) {$O$};
\end{tikzpicture}
\end{figure}


A particle $P$ of mass $m$ is attached to one end of a light elastic string of natural length $5a$ and modulus of elasticity $3mg$. The other end of the string is fixed to a point $O$ on a smooth plane, which is inclined at an angle $\alpha$ to the horizontal, where $\sin\alpha=\frac{3}{5}$. The string lies along a line of greatest slope of the plane and $P$ rests in equilibrium on the plane at a point $A$ below $O$, as shown in the diagram above. \\
At time $t=0$, $P$ is projected from $A$ with speed $2\sqrt{\frac{3ag}{5}}$ up the line of greatest slope towards $O$. At time $t$, the displacement of $P$ from $A$ is $x$, where the positive direction of $x$ is up the plane. \\
\begin{enumerate}[label=\textbf{(\alph*)}]
\item Show that, while the string is taut, $x$ satisfies the differential equation
\[
\frac{\mathrm{d}^2x}{\mathrm{d}t^2}+\frac{3gx}{5a}=0
\] \marks[-30pt]{6}
\item Find, in terms of $a$ and $g$, the time at which the string first becomes slack. \marks{5} \\
\item Find, in terms of $a$, the distance from $A$ of the highest point on the plane reached by $P$. \marks{3} \\
\end{enumerate}

\newpage

% ---- Question 9 ----
\questionitem
% [Source: _QBT__Simple_Harmonic_Motion, Q9]

\begin{figure}[H]
    \centering
\begin{tikzpicture}[scale=1]
\def\AB{10}
\def\AE{3}
\pgfmathsetmacro{\CE}{2/3}
\def\ext{0.6}
\def\tick{0.15}
\def\ctick{0.1}
\def\dimLow{0.7}
\def\dimHigh{1.5}
\def\prad{2.5pt}

\coordinate (A) at (0,0);
\coordinate (B) at (\AB,0);
\coordinate (E) at (\AE,0);
\coordinate (C) at ($(E)+(-\CE,0)$);

\draw[thick] ($(A)+(-\ext,0)$) -- ($(B)+(\ext,0)$);

\draw[thick] ($(A)+(0,-\tick)$) -- ($(A)+(0,\tick)$);
\draw[thick] ($(B)+(0,-\tick)$) -- ($(B)+(0,\tick)$);
\draw[thick] ($(C)+(0,-\ctick)$) -- ($(C)+(0,\ctick)$);

\fill[black] (E) circle (\prad);

\draw[{Stealth[length=2mm]}-{Stealth[length=2mm]}] ($(C)+(0,\dimLow)$) -- ($(E)+(0,\dimLow)$)
  node[midway, above] {$\frac{2}{3}L$};
\draw[{Stealth[length=2.5mm]}-{Stealth[length=2.5mm]}] ($(A)+(0,\dimHigh)$) -- ($(B)+(0,\dimHigh)$)
  node[midway, above] {$10L$};

\node[below=3pt] at (A) {$A$};
\node[below=3pt] at (B) {$B$};
\node[below=3pt] at (C) {$C$};
\node[below=3pt] at (E) {$E$};
\node[above right=3pt] at (E) {$P$};
\end{tikzpicture}
\end{figure}


A particle $P$ of mass $m$ lies on a smooth horizontal table. \\
One end of a light elastic string, of natural length $2L$ and modulus of elasticity $3mg$, is attached to $P$, and the other end is attached to a fixed point $A$ on the table. \\
One end of a second light elastic string, of natural length $4L$ and modulus of elasticity $2mg$, is also attached to $P$, and the other end is attached to a fixed point $B$ on the table, where $AB=10L$. \\
The particle $P$ rests in equilibrium at the point $E$, where $AEB$ is a straight line, as shown in the diagram above. \\
\begin{enumerate}[label=\textbf{(\alph*)}]
\item Show that $AE=3L$. \marks{2} \\
\end{enumerate}
The particle $P$ is now held at the point $C$ on $AE$, where $CE=\frac{2}{3}L$, and released from rest. \\
\begin{enumerate}[label=\textbf{(\alph*)}, resume]
\item Show that $P$ moves with simple harmonic motion with centre $E$ and period $\pi\sqrt{\frac{2L}{g}}$. \marks{6} \\
\item Find, in terms of $L$ and $g$, the speed of $P$ at the instant when $AP=\frac{5}{2}L$. \marks{3} \\
\item Find, in terms of $L$ and $g$, the exact length of time in each complete oscillation for which the tension in the string $AP$ is greater than $2mg$. \marks{6} \\
\end{enumerate}

\newpage

% ---- Question 10 ----
\questionitem
% [Source: _QBT__Simple_Harmonic_Motion, Q10]

A particle $P$ of mass $m$ is attached to one end of a light elastic string of natural length $l$ and modulus of elasticity $mg$. The other end of the string is attached to a fixed point $O$. \\
Particle $P$ is projected vertically downwards from $O$ with speed $\sqrt{gl}$. \\
\begin{enumerate}[label=\textbf{(\alph*)}]
\item Show that the greatest length of the string in the subsequent motion is $4l$. \marks{5} \\
\end{enumerate}
The string first becomes taut at time $t=0$. At time $t \ge 0$, the length of the string is $l+x$, where $x$ is the extension in the string. \\
\begin{enumerate}[label=\textbf{(\alph*)},resume]
\item Show that, while the string is taut, $x$ satisfies the differential equation
\[
\ddot{x}+\omega^2 x=g, \qquad \text{where } \omega^2=\dfrac{g}{l}
\]
\marks[-30pt]{3}
\item By using the substitution $x=y+\frac{g}{\omega^2}$, solve the differential equation and hence find, in terms of $l$ and $g$, the exact length of time for which the string is taut before it first becomes slack. \marks{12} \\
\end{enumerate}

\end{enumerate}
\end{document}
